% !TeX root = ../../金融经济学_熊和平_习题.tex
% 源：微信图片_20260928151500_3593_382.jpg（教材印张 311）

\zhchapter{异质性投资者}

\begin{tcolorbox}[summary,title=本章小结]
在行为金融中所重现的一些心理特征都可能形成特定的异质性；本章的异质性属于一类最基本的异质。即使是上述基本的异质性也会对资产定价和最终的资源配置产生很大的影响。

学术界针对经济人的异质性通常有三大假设：一是市场平衡假设，认为众多异质性经济人相互作用最终可以将诸多个性“平衡”掉，从而形成貌似只有一个经济人的情形；第二是市场选择假设，借助生物中的“物竞天择、自然选择”思想，认为市场将会淘汰“错误的”经济人，最终使聪明的经济人留下，从而使得资本市场被聪明的经济人所主导；第三是市场学习假设，认为经济人利用可获得的信息更新自己对市场的认知，从而形成正确的预期并最终使市场趋于理想化的市场。三种假设的模型化和实证分析都推动金融经济学理论的发展。

市场平衡假设强调所有经济人的“集体”作用，不同经济人共同作用形成了资本市场的加总理论，异质禀赋下的加总形成加总性质，CRRA 偏好的加总则形成一类简单的、理性化的情形。
\end{tcolorbox}

\begin{tcolorbox}[knowledge,title=重要概念]
异质性投资者，异质禀赋，异质偏好，异质信念，时间偏好，资本市场加总，市场选择，市场平衡，投资者学习，贝叶斯更新，加总性质。
\end{tcolorbox}

% !TeX root = ../../金融经济学_熊和平_习题.tex

\xhpthoughts
\begin{problem}
\problemno{1} 投资者的异质性一般包括哪几个方面？
\end{problem}
\begin{problem}
\problemno{2} 什么是投资者的加总性质？
\end{problem}
\begin{problem}
\problemno{3} 关于异质性经济人有哪三大假设？
\end{problem}
\begin{problem}
\problemno{4} 加总性质与特定效用函数下的加总有哪些区别与联系？
\end{problem}
\begin{problem}
\problemno{5} 如何理解理性学习与贝叶斯更新？
\end{problem}

\xhpexercises
\begin{problem}
\problemno{1} 在例 14.14 中取 $B=1$，$A=0$ 时，效用函数为对数偏好。请在对数偏好下对经济人进行加总。
\end{problem}
\begin{problem}
\problemno{2} 假设经济中两个经济人的时间偏好分别为 $0.5$ 和 $1$，他们都具有 CRRA 效用函数，他们的相对风险厌恶系数为 $\gamma=1$，分别求例 14.1 中三种异质禀赋下的状态价格和资源配置。
\end{problem}
\begin{problem}
\problemno{3} 在上题的条件下进一步假设两个经济人具有异质性偏好，相对风险厌恶系数分别为 $1$ 和 $0.1$。在此条件下，分别求例 14.1 中三种异质禀赋下的状态价格和资源配置。
\end{problem}
\begin{problem}
\problemno{4} 假设经济中两个经济人的时间偏好均为 $1$，他们都具有 CRRA 效用函数，他们的相对风险厌恶系数为 $1$ 和 $0.5$，分别求例 14.1 中三种异质禀赋下的状态价格和资源配置。
\end{problem}

